\documentclass[11pt]{article}

\usepackage[margin=1in]{geometry}
\usepackage{amsmath,amssymb,booktabs,graphicx,hyperref}

\title{TEMPEST Model Bake-Off}
\author{TEMPEST Model Bake-Off contributors}
\date{\today}

\begin{document}
\maketitle

\begin{abstract}
This report will document a mathematical and computational closure program
spanning particle-to-kinetic, wave-to-kinetic, kinetic-to-fluid, and
application-scale reductions motivated by TEMPEST.  It will compare classical
Markov closures, explicit cumulant and memory models, physics-informed neural
networks, neural operators, physics-informed neural operators, learned
closures, structured learned dynamics, and action-conditioned world models.
The first application profile is nonlocal magnetized heat transport for
magnetic target fusion.
\end{abstract}

\setcounter{tocdepth}{1}
\tableofcontents

\section{Introduction}
The organizing question is not which architecture wins in the abstract, but
what information is required to close a reduced model reliably and how that
requirement changes across asymptotic and breakdown regimes.

We built an evaluation harness for computational closure methods: a fixed
suite of test problems and forcing regimes, a common interface through which
any closure model---analytic, learned, or hybrid---is invoked, and a
standardized battery of accuracy, stability, and cost metrics computed
uniformly across candidates. The harness separates the science (the closures)
from the scaffolding (data handling, integration, scoring), so that adding a
new method requires implementing one interface, and every reported comparison
is reproducible and fair by construction.

\section{Mathematical closure framework}
This section will define scale transitions and exact hierarchy defects,
connected correlations and cumulants, interaction histories and finite memory,
and the decomposition of numerical, statistical, asymptotic, closure, and
learning error.

For a representative hierarchy,
\begin{equation}
  f_2=f_1f_1+g_2,
  \qquad
  \partial_t f_1=\mathcal L f_1+\mathcal C[f_1]
  +\mathcal R[g_2,g_3,\ldots],
\end{equation}
where the closure problem is to approximate the residual $\mathcal R$ from a
declared information set while preserving the required mathematical structure.

\section{Canonical scale-limit benchmarks}
The canonical track is motivated by the long-time scale-limit program of
Zaher Hani and collaborators. It contains hard-particle dynamics to
Boltzmann, nonlinear wave dynamics to the wave kinetic equation, and
Boltzmann dynamics to Euler and Navier--Stokes--Fourier. These are not
ready-made TEMPEST closures. They are benchmark regimes in which the
microscopic and reduced equations, asymptotic coordinates, expected limit,
and discarded information are all explicit.

This makes them the validation layer for the computational evaluation
harness. The particle track exposes connected cumulants, recollisions, and
collision ancestry; the wave track exposes spectral cumulants and resonant
interaction histories; and the kinetic-to-fluid track exposes
non-Maxwellian moments and constitutive remainders. Each benchmark will test
the trusted asymptotic regime, controlled breakdown outside its assumptions,
long-time behavior, and separation of numerical, ensemble, asymptotic, and
closure-model errors. TEMPEST applications retain this diagnostic protocol
but replace the equations, states, constraints, observables, and coefficients.

\section{TEMPEST application profiles}
The flagship profile will study magnetic-target-fusion nonlocal heat
transport, including an explicit canonical-to-application translation. A
dusty-plasma transfer profile remains a candidate second application.

\section{Methods and information budgets}
Methods will include classical, cumulant, and memory closures; PINNs, neural
operators, and PINOs; learned closures and structured dynamics; and
action-conditioned world models. Privileged information and matched
comparisons will be reported explicitly.

\section{Experimental protocol}
The protocol will cover reference-solver verification, ensemble and estimator
uncertainty, a-priori closure tests, a-posteriori embedded-solver tests, and
transfer and failure-boundary experiments.

\section{Results}
Results will report recovery of classical limits; cumulant and memory
diagnostics; predictive and embedded closure performance; MTF
application-profile performance; and compute and information efficiency.

\section{Discussion}
The discussion will ask which closure information is necessary, what transfers
from canonical systems to MTF, and where negative results limit
generalization.

\section{Reproducibility statement}
The reproducibility record will link data and result manifests, code,
configurations, environments, claim levels, and retained failures.

\end{document}
